\input{preamble}

%----------------------------------------------------------------------------------
%            content
%----------------------------------------------------------------------------------
\begin{document}
%\begin{CJK*}{UTF8}{gbsn}                          % to typeset your resume in Chinese using CJK
%-----       resume       ---------------------------------------------------------
\makecvtitle
\vspace{-25pt}
\section{Research Interests}
\cvitem{}{High Performance Computing, Compilers, Debuggers, and AI for SE}

\section{Education}
\cventry{2024 - present}{Ph.D.}{The University of Texas at Austin}{}{\textit{GPA: 3.9}}{Major in Software Engineering and Systems, advised by Milos Gligoric.}
\cventry{2020 - 2024}{B.S.}{Belarusian State University of Informatics and Radioelectronics}{}{\textit{GPA: 3.7}}{Major in Computer Science and Network Engineering}

% \section{Master thesis}
% \cvitem{title}{\emph{Title}}
% \cvitem{supervisors}{Supervisors}
% \cvitem{description}{Short thesis abstract}

\section{Experience}
\cventry{2024 - present}{Graduate Research Assistant}{University of Texas at Austin}{}{}{Software Validation, Compilers and High Performance Computing
      \vspace{2pt}
      \begin{itemize}
            \item Implemented a debugger for HPC eDSL with $\leq2.33\times$ overhead compared to \pdb, \Python{'s} native debugger;
            \item Developed static and dynamic analyzer tools for the existing programming language on top of its infrastructure;
            \item Main developer and maintainer of \pykokkos, contributor of Kokkos;
            \item Research ML for HPC optimizations, HPC eDSLs, and Python-to-C++ interoperability.
            \item Research compiler optimization techniques, source code processing, and code testing tools;
      \end{itemize}}

\cventry{Summer 2026}{PhD Intern}{Oak Ridge National Laboratory}{\href{https://github.com/arborx/pyArborX}{\underline{pyArborX}}}{}{
      \begin{itemize}
            \item Designed and built a JIT compiler translating Python into C++ for ArborX, an HPC spatial-search library.
            \item Deployed a custom AST parser and py2many-based code generator for ArborX's template-heavy C++ APIs.
            \item Implemented a pybind11 module generator that enables JIT compilation and delayed template specialization.
            \item Validated framework by implementing ML clustering algorithms and deploying CIs for correctness checking.
      \end{itemize}
}

\cventry{Spring 2024}{Research Assistant}{BSUIR}{}{}{Led the research on wireless communication technology
      \vspace{2pt}
      \begin{itemize}
            % \item Researched the wireless networks, like Bluetooth and WiFi, for interference;
            \item Conducted in-depth research on the topic of harmonic oscillations and total harmonic distortion between signals in wireless networks;
                  % \item Developed the formal method of defining interference between devices;
            \item Developed and tested the PCB for jamming Bluetooth and WiFi devices in the shielded lab;
                  % \item Tested the developed PCB and found corner cases of the technologies mentioned above;
            \item Contributed to 3 different research projects over 5 months.
      \end{itemize}}

\cventry{Summer 2023}{Software Developer Intern}{AGAT - Control Systems}{}{}{Contributed to multiple network tools for the company network.
      \vspace{2pt}
      \begin{itemize}
            \item Developed network controlling and monitoring software for Cisco and Dell equipment.
            \item Implemented and designed network analysis tools for the company's local network.
            \item Collaborated with the team to deploy an improved network management system.
      \end{itemize}}

\section{Publications}
\cvitem{SC '26}{\paperentry{Ivan Grigorik, Gabriel Kosmacher, George Biros, Milos Gligoric}{Interactive Debugger for Performance Portable Python HPC Kernels}}
\cvitem{SciConf '24}{\paperentry{Ivan Grigorik, Diana Kupriyanova}{Active jammer stations based on direct influence noises}}

\section{Skills}
\cvitem{Languages}{C, C++,  Python, Rust, ASM, Java, Bash, SQL;}
\cvitem{Tools}{GCC, LLVM, Make \& CMake, CUDA, GDB, Valgrind, Docker, STL;}
\cvitem{Domains}{Compilers, HPC, Static and Dynamic Analysis, Verification of Software;}
\cvitem{Concepts}{System and Software Engineering, Program Testing, Sanitizing and Profiling, Design Patterns, Distributed Computing, Parallel Algorithms, Multithreading, Computer Architecture, Networking;}
\newpage

\section{Services and event participation}
\cvitem{2026}{CGO '27 -- artifact evaluator.\newline
      CGO '27 -- sub-reviewer.\newline
      OOPSLA '26 -- sub-reviewer.\newline
      UT '26 -- PhD student participant in faculty hiring process.}
\cvitem{2025}{CGO '25 -- sub-reviewer}
\subsection{Presentations}
% \cvitem {TODO: add SC26 presentation here + pkdb paper with the ref to the SC journal}

\cvitem{2026}{SC '26: Interactive Debugger for Performance Portable Python HPC kernels}
\cvitem{2026}{AMD iMAGiNE Symposium.}
\cvitem{2026}{Predictive Science Academic Alliance Program (PSAAP) presentation: \pykokkos ecosystem and development cycle. }
\cvitem{2026}{High Performance Software Foundation (HPSF) presentation: Pythonizing \pykokkos.}
\cvitem{2025}{AMD iMAGiNE Symposium.}
\cvitem{2024 - 2026}{Co-organized the Joint UT-Cornell Software Engineering Seminar Series.}

\section{Projects}
\cvitem{xbash}{A fork of bash with static and dynamic analysis tools and extended builtins.}
\cvitem{cuRTX}{Open-source ray tracing tool written in CUDA and C++, inspired by RTX-in-one-weekend book.}
\cvitem{GameOfLife}{Pet-project, Conway's Game of Life written in C++ and SFML.}
\cvitem{HvC}{Fork of HIPIFY that compares heterogeneous computing platforms (CUDA, HIP, etc.).}
\cvitem{\pykokkos}{Framework for writing high-performance Python code similar to Numba. *Current maintainer.}

% \section{Relevant coursework}

% \cvlistdoubleitem{Programming languages and computer design}{System software}
% \cvlistdoubleitem {Computer Architecture}{HPC architecture}
% \cvlistdoubleitem {Software Networks}{Hardware Systems}
% \cvlistdoubleitem {Compilers}{Computer Graphics}

\section{Honors and awards}
\cvitem{2024}{Graduate School Fellowship, University of Texas at Austin.}
\cvitem{2024}{Outstanding paper of BSUIR 60th SciConf24 \newline Paper: Active jammer stations based on direct influence noises.}

%   \begin{itemize}
%   \item Sub-achievement (a);
%   \item Sub-achievement (b), with sub-sub-achievements (don't do this!);
%     \begin{itemize}
%     \item Sub-sub-achievement i;
%     \item Sub-sub-achievement ii;
%     \item Sub-sub-achievement iii;
%     \end{itemize}
%   \item Sub-achievement (c);
%   \end{itemize}
% \item Achievement 3
% \item Achievement 4
% \cventry{year--year}{Job title}{Employer}{City}{}{Description line 1\newline{}Description line 2\newline{}Description line 3}
% \subsection{Miscellaneous}
% \cventry{year--year}{Job title}{Employer}{City}{}{Description}

% \section{Languages}
% \cvitemwithcomment{Language 1}{Skill level}{Comment}
% \cvitemwithcomment{Language 2}{Skill level}{Comment}
% \cvitemwithcomment{Language 3}{Skill level}{Comment}
% \cvitemwithcomment{Language 4}{Skill level}{Comment}

% \section{Computer skills}
% \cvdoubleitem{category 1}{XXX, YYY, ZZZ}{category 4}{XXX, YYY, ZZZ}
% \cvdoubleitem{category 2}{XXX, YYY, ZZZ}{category 5}{XXX, YYY, ZZZ}
% \cvdoubleitem{category 3}{XXX, YYY, ZZZ}{category 6}{XXX, YYY, ZZZ}
%% Skill matrix as an alternative to rate one's skills, computer or other. 

%% Adjusts width of skill matrix columns. 
%% Usage \setcvskillcolumns[<width>][<factor>][<exp_width>]
%% <width>, <exp_width> should be lengths smaller than \textwidth, <factor> needs to be between 0 and 1.
%% Examples:
% \setcvskillcolumns[5em][][]%    adjust first column. Same as \setcvskillcolumns[5em]
% \setcvskillcolumns[][0.45][]%   adjust third (skill) column. Same as \setcvskillcolumns[][0.45]
% \setcvskillcolumns[][][\widthof{``Year''}]%     adjust fourth (years) column.
% \setcvskillcolumns[][0.45][\widthof{``Year''}]%
% \setcvskillcolumns[\widthof{``Languag''}][0.48][]
% \setcvskillcolumns[\widthof{``Languag''}]%

%% Adjusts width of legend columns. Usage \setcvskilllegendcolumns[<width>][<factor>]
%% <factor> needs to be between 0 and 1. <width> should be a length smaller than \textwidth
%% Examples:
% \setcvskilllegendcolumns[][0.45]
% \setcvskilllegendcolumns[\widthof{``Legend''}][0.45]
% \setcvskilllegendcolumns[0ex][0.46]% this is usefull for the banking style

%% Add a legend if you are using \cvskill{<1-5>} command or \cvskillentry
%% Usage \cvskilllegend[*][<post_padding>][<first_level>][<second_level>][<third_level>][<fourth_level>][<fifth_level>]{<name>}
% \cvskilllegend % insert default legend without lines
% \cvskilllegend*[1em]{}% adjust post spacing
% \cvskilllegend*{Legend}%  Alternatively add a description string
%% adjust the legend entries for other languages, here German
% \cvskilllegend[0.2em][Grundkenntnisse][Grundkenntnisse und eigene Erfahrung in Projekten][Umfangreiche Erfahrung in Projekten][Vertiefte Expertenkenntnisse][Experte\,/\,Spezialist]{Legende}

%% Alternative legend style with the first three skill levels in one column
%% Usage \cvskillplainlegend[*][<post_padding>][<first_level>][<second_level>][<third_level>][<fourth_level>][<fifth_level>]{<name>}
% \setcvskilllegendcolumns[][0.6]%  works for classic, casual, banking
% \setcvskilllegendcolumns[][0.55]%  works better for oldstyle and fancy
% \cvskillplainlegend{}
% \cvskillplainlegend[0.2em][Grundkenntnisse][Grundkenntnisse und eigene Erfahrung in Projekten][Umfangreiche Erfahrung in Projekten][Vertiefte Expertenkenntnisse][Experte/Guru]{Legende}

%% Add a head of the skill matrix table with descriptions.
%% Usage \cvskillhead[<post_padding>][<Level>][<Skill>][<Years>][<Comment>]%
% \cvskillhead[-0.1em]%   this inserts the standard legend in english and adjust padding
%% Adjust head of the skill matrix for other languages
% \cvskillhead[0.25em][Level][F\"ahigkeit][Jahre][Bemerkung]

%% \cvskillentry[*][<post_padding>]{<skill_cathegory>}{<0-5>}{<skill_name>}{<years_of_experience>}{<comment>}% 
%% Example usages:
% \cvskillentry*{Language:}{3}{Python}{2}{I'm so experienced in Python and have realised a million projects. At least.}
% \cvskillentry{}{2}{Lilypond}{14}{So much sheet music! Man, I'm the best!}
% \cvskillentry{}{3}{\LaTeX}{14}{Clearly I rock at \LaTeX}
% \cvskillentry*{OS:}{3}{Linux}{2}{I only use Archlinux btw}% notice the use of the starred command and the optional 
% \cvskillentry*[1em]{Methods}{4}{SCRUM}{8}{SCRUM master for 5 years}
%% \cvskill{<0-5>} command
% \cvitem{\textbackslash{cvskill}:}{Skills can be visually expressed by the \textbackslash{cvskill} command, e.g. \cvskill{2}}

% \section{Soft skills}
% \cvitem{hobby 1}{Description}
% \cvitem{hobby 2}{Description}
% \cvitem{hobby 3}{Description}

% \section{Interests}
% \cvitem{ }{Decription 1}
% \cvitem{ }{Decription 2}
% \cvitem{ }{Decription 3}

% \section{Extra 1}
% \cvlistitem{Item 1}
% \cvlistitem{Item 2}
% \cvlistitem{Item 3}

% \section{Extra 2}
% \cvlistdoubleitem{Item 1}{Item 4}
% \cvlistdoubleitem{Item 2}{Item 5\cite{book2}}
% \cvlistdoubleitem{Item 3}{Item 6. Like item 3 in the single column list before, this item is particularly long to wrap over several lines.}

% \section{References}
% \begin{cvcolumns}
%   \cvcolumn{Category 1}{\begin{itemize}\item Person 1\item Person 2\item Person 3\end{itemize}}
%   \cvcolumn{Category 2}{Amongst others:\begin{itemize}\item Person 1, and\item Person 2\end{itemize}(more upon request)}
%   \cvcolumn[0.5]{All the rest \& some more}{\textit{That} person, and \textbf{those} also (all available upon request).}
% \end{cvcolumns}

% % Publications from a BibTeX file without multibib
% %  for numerical labels: \renewcommand{\bibliographyitemlabel}{\@biblabel{\arabic{enumiv}}}% CONSIDER MERGING WITH PREAMBLE PART
% %  to redefine the heading string ("Publications"): \renewcommand{\refname}{Articles}
% \nocite{*}
% \bibliographystyle{plain}
% \bibliography{publications}                        % 'publications' is the name of a BibTeX file

% Publications from a BibTeX file using the multibib package
% \nocitebook{book1,book2}
% \bibliographystylebook{plain}
% \bibliographybook{publications}                   % 'publications' is the name of a BibTeX file
% \nocitemisc{misc1,misc2,misc3}
% \bibliographystylemisc{plain}
% 'publications' is the name of a BibTeX file

%%%%%%%%%%%%%%%%%%%%%%%%%%%%%%% Cover letter
\clearpage
% %-----       letter       ---------------------------------------------------------
% recipient data
\recipient{\Fix{change recipient}}{}
\date{\today}
\opening{Dear \Fix{<WHO>}}
\closing{Sincerely,}
\makelettertitle

Dear \Fix{INPUT},

\Fix{Input text here}

\makeletterclosing

%\clearpage\end{CJK*}                              % if you are typesetting your resume in Chinese using CJK; the \clearpage is required for fancyhdr to work correctly with CJK, though it kills the page numbering by making \lastpage undefined
\end{document}

%% end of file `template.tex'.

